\chapter{Belajar dari Ganjaran: Reinforcement Learning}
\label{ch:rl}

Reinforcement learning datang kepada saya dalam dua tahap. Pada musim panas 2025, kuliah deep
reinforcement learning langsung membawa kami ke imitation learning, policy gradient, dan credit
assignment, dengan proyek latihan berupa agen DQN, agen PPO, dan sebuah tantangan akhir. Pada
musim dingin berikutnya, kuliah reinforcement learning mengisi fondasinya secara rapi: proses
keputusan Markov, dynamic programming, metode Monte Carlo, dan temporal difference. Bab ini
menyusunnya dalam urutan yang lebih logis, fondasi dulu, lalu deep RL.

Dua tipe belajar di bab-bab sebelumnya selalu punya data yang sudah tersedia. Reinforcement learning
berbeda. Agen harus mengumpulkan datanya sendiri dengan bertindak, dan umpan baliknya bukan jawaban
yang benar, melainkan angka ganjaran yang sering datang terlambat. Seekor anjing yang belajar duduk
tidak diberi tahu gerakan otot apa yang harus dilakukan. Ia mencoba sesuatu, dan kadang mendapat
camilan. Tugasnya adalah mencari tahu tindakan mana yang membuat camilan itu datang.

\section{Proses keputusan Markov}

Kerangka formalnya adalah \istilah{Markov decision process} (MDP): himpunan keadaan $\mathcal S$,
himpunan aksi $\mathcal A$, dinamika $p(s',r\mid s,a)$ yang memberi peluang keadaan berikutnya dan
ganjarannya, serta faktor diskon $\gamma\in[0,1)$ \citep{sutton2018}. Di setiap langkah, agen melihat
keadaan $s_t$, memilih aksi $a_t$ menurut \istilah{policy} $\pi(a\mid s)$, lalu menerima ganjaran
$r_{t+1}$ dan keadaan baru. Tujuannya memaksimalkan \istilah{return}, jumlah ganjaran yang didiskon,
\begin{equation}
G_t=r_{t+1}+\gamma r_{t+2}+\gamma^2r_{t+3}+\cdots=\sum_{k=0}^\infty\gamma^kr_{t+k+1}.
\end{equation}
Faktor diskon membuat ganjaran yang jauh di masa depan bernilai lebih kecil, dan secara matematis
menjaga jumlahnya tetap terhingga.

Seberapa baik berada di sebuah keadaan diukur oleh \istilah{value function} $V^\pi(s)=\E_\pi[G_t\mid s_t=s]$,
dan seberapa baik mengambil sebuah aksi di sebuah keadaan oleh \istilah{action-value function}
$Q^\pi(s,a)=\E_\pi[G_t\mid s_t=s,a_t=a]$. Kedua fungsi memenuhi persamaan rekursif yang menjadi jantung
seluruh bidang ini, yaitu persamaan Bellman:
\begin{equation}
V^\pi(s)=\sum_a\pi(a\mid s)\sum_{s',r}p(s',r\mid s,a)\big[r+\gamma V^\pi(s')\big].
\end{equation}
Nilai sebuah keadaan sama dengan ganjaran langsung ditambah nilai keadaan berikutnya yang didiskon,
dirata-ratakan atas semua yang mungkin terjadi. Untuk policy terbaik, rata-rata atas aksi diganti
dengan maksimum, dan kita mendapat persamaan Bellman optimal,
$V^*(s)=\max_a\sum_{s',r}p(s',r\mid s,a)[r+\gamma V^*(s')]$.

Persamaan ini turun langsung dari definisi return. Karena $G_t=r_{t+1}+\gamma G_{t+1}$, ekspektasinya
dari keadaan $s$ adalah $\E[r_{t+1}+\gamma G_{t+1}\mid s_t=s]$. Ekspektasi ini dihitung bertahap: pilih
aksi menurut $\pi$, dapatkan ganjaran dan keadaan berikutnya menurut $p$, lalu ganti
$\E[G_{t+1}\mid s_{t+1}=s']$ dengan $V^\pi(s')$ berkat sifat Markov, karena masa depan setelah $s'$ tidak
bergantung pada bagaimana kita sampai di sana.

Contoh paling kecil sudah memperlihatkan maknanya. Bayangkan satu keadaan dengan satu aksi yang
selalu memberi ganjaran satu dan kembali ke keadaan yang sama. Persamaan Bellman menjadi
$V=1+\gamma V$, sehingga $V=1/(1-\gamma)$. Dengan $\gamma=0{,}9$, nilainya 10. Faktor diskon bisa
dibaca sebagai cakrawala efektif: dengan $\gamma=0{,}9$, agen kira-kira peduli pada sepuluh langkah ke
depan.

\section{Dynamic programming}

Kalau dinamika MDP diketahui, persamaan Bellman bisa diselesaikan dengan dynamic programming
\citep{bellman1957}. \istilah{Policy evaluation} menghitung $V^\pi$ dengan menerapkan persamaan Bellman
berulang-ulang sebagai aturan update sampai nilainya tidak berubah lagi. Karena faktor diskon
membuat update ini kontraktif, iterasinya pasti konvergen. \istilah{Policy improvement} lalu membuat
policy baru yang serakah terhadap $V^\pi$, memilih di setiap keadaan aksi dengan nilai harapan
tertinggi. Teorema perbaikan policy menjamin policy baru ini tidak lebih buruk dari yang lama.
Bergantian antara keduanya disebut \istilah{policy iteration}, dan ia berakhir di policy optimal
dalam jumlah langkah terhingga untuk MDP terhingga. \istilah{Value iteration} memadatkan keduanya
dengan langsung menerapkan persamaan Bellman optimal sebagai update.

Kuliah memakai dunia grid kecil sebagai contoh. Agen bergerak di petak-petak, setiap langkah memberi
ganjaran $-1$, dan permainan berakhir di petak sudut. Setelah policy evaluation untuk policy acak,
nilai setiap petak mencerminkan rata-rata jumlah langkah yang dibutuhkan untuk mencapai sudut
dengan berjalan acak. Policy yang serakah terhadap nilai-nilai itu sudah langsung menunjuk jalur
terpendek.

Value iteration bisa diikuti dengan tangan pada lorong empat petak (\cref{fig:lorong}). Agen bisa melangkah ke kiri atau
ke kanan, setiap langkah memberi $-1$, petak paling kanan adalah tujuan dengan nilai nol, dan $\gamma=1$. Mulai dari
$V=0$ di mana-mana. Setelah satu sapuan, setiap petak selain tujuan bernilai $-1$, karena satu langkah apa pun memberi
$-1$ dan nilai petak tetangga masih nol. Setelah sapuan kedua, petak tepat di sebelah tujuan tetap $-1$ (melangkah ke
tujuan), sedangkan petak lain menjadi $-2$. Setelah sapuan ketiga, nilainya $(-3,-2,-1,0)$, dan sapuan berikutnya tidak
mengubah apa pun. Nilai setiap petak adalah minus jarak ke tujuan, dan policy serakah selalu melangkah ke kanan.
Informasi menyebar mundur dari tujuan satu petak per sapuan, seperti gelombang. Pada MDP yang lebih besar, gelombang
yang sama juga merambat, hanya lebih lama.

\begin{figure}[ht]
\centering
\begin{tikzpicture}[x=1.6cm,y=1cm]
  \foreach \i/\v in {0/$-3$,1/$-2$,2/$-1$}{
    \draw[abu] (\i,0) rectangle (\i+1,1);
    \node at (\i+0.5,0.5) {\v};
    \draw[panah,oranye] (\i+0.62,0.15) -- (\i+0.95,0.15);
  }
  \fill[hijau!25] (3,0) rectangle (4,1);
  \draw[abu] (3,0) rectangle (4,1);
  \node at (3.5,0.5) {$0$};
  \node[label,below] at (3.5,0) {tujuan};
  \node[label,above] at (2,1) {setiap langkah: $-1$};
\end{tikzpicture}
\caption{Lorong empat petak setelah value iteration konvergen. Nilai setiap petak adalah minus jaraknya ke tujuan, dan
policy serakah (panah) selalu melangkah ke kanan.}
\label{fig:lorong}
\end{figure}

\section{Belajar dari pengalaman}

Di dunia nyata, dinamika jarang diketahui. Agen harus belajar dari sampel.

Metode \istilah{Monte Carlo} menunggu sampai satu episode selesai, lalu memakai return yang
benar-benar diperoleh sebagai perkiraan nilai setiap keadaan yang dikunjungi. Perkiraan ini tak
bias, tetapi variansnya besar, karena satu return bergantung pada semua keputusan acak sampai akhir
episode. Kuliah memakai blackjack sebagai contoh: setelah ratusan ribu permainan acak dengan sebuah policy sederhana, rata-rata
return dari setiap keadaan kartu menggambar permukaan nilai yang mulus, tanpa pernah menuliskan peluang kartu. Ada dua
varian. \emph{First-visit} hanya memakai return sejak kunjungan pertama ke sebuah keadaan dalam satu episode, \emph{every-visit}
memakai setiap kunjungan. Untuk episode $s_1,s_2,s_1,\text{selesai}$ dengan ganjaran $1$, $0$, dan $2$ dan $\gamma=1$, return sejak
kunjungan pertama ke $s_1$ adalah $3$ dan sejak kunjungan kedua $2$. First-visit mencatat $3$, every-visit mencatat $3$ dan $2$.
Keduanya konvergen ke nilai yang benar. Rata-rata ini bisa diperbarui bertahap, $V\leftarrow V+\tfrac1n(G-V)$, bentuk yang
langsung menjadi model untuk TD di bawah dengan $\tfrac1n$ diganti laju belajar tetap.

Metode \istilah{temporal difference} (TD) tidak menunggu. Setelah satu langkah, ia memperbarui
perkiraan dengan memakai perkiraannya sendiri tentang keadaan berikutnya,
\begin{equation}
V(s_t)\leftarrow V(s_t)+\alpha\big[r_{t+1}+\gamma V(s_{t+1})-V(s_t)\big].
\end{equation}
Selisih di dalam kurung disebut \istilah{TD error}. Pendekatan ini, belajar tebakan dari tebakan,
disebut \istilah{bootstrapping}. Variansnya kecil karena hanya satu langkah acak yang terlibat, tetapi
ada bias selama perkiraan $V(s_{t+1})$ belum benar. Di antara keduanya terbentang metode $n$-langkah
dan eligibility traces yang mencampur keduanya.

Bayangkan memperkirakan waktu perjalanan pulang dari kantor. Metode Monte Carlo menunggu sampai Anda
tiba di rumah, lalu mencatat total waktunya. Metode TD memperbarui perkiraan di setiap persimpangan:
``jalanan ternyata macet, perkiraan saya untuk sisa perjalanan harus naik'', tanpa menunggu tiba.

Untuk mengendalikan, bukan hanya menilai, kita mempelajari $Q$. SARSA memperbarui $Q(s,a)$ dengan aksi
berikutnya yang benar-benar diambil, sehingga ia belajar nilai policy yang sedang dijalankan.
Q-learning \citep{watkins1992} memakai aksi terbaik di keadaan berikutnya,
\begin{equation}
Q(s_t,a_t)\leftarrow Q(s_t,a_t)+\alpha\big[r_{t+1}+\gamma\max_{a'}Q(s_{t+1},a')-Q(s_t,a_t)\big],
\end{equation}
sehingga ia belajar nilai policy optimal sambil menjalankan policy lain yang lebih eksploratif.
Perbedaan ini disebut on-policy dan off-policy.

Satu update dengan angka: misalkan $Q(s,a)=2$, agen menerima $r=1$, dan di keadaan berikutnya nilai aksi terbaik
$\max_{a'}Q(s',a')=5$. Dengan $\gamma=0{,}9$ dan $\alpha=0{,}1$, targetnya $1+0{,}9\cdot5=5{,}5$, TD error-nya $3{,}5$, dan
$Q(s,a)$ naik menjadi $2{,}35$. Langkah kecil ke arah target, persis seperti SGD dengan target yang diambil dari
perkiraan sendiri.

Perbedaan on-policy dan off-policy terlihat jelas pada contoh klasik \citet{sutton2018}, berjalan di tepi tebing.
Jalur terpendek menyusuri tepi tebing, dan setiap langkah yang jatuh ke tebing memberi hukuman besar. Dengan eksplorasi
$\epsilon$-greedy, sesekali agen melangkah acak. Q-learning mempelajari nilai policy optimal yang tidak pernah salah
langkah, sehingga ia memilih jalur di tepi tebing, lalu sesekali jatuh karena eksplorasinya. SARSA mempelajari nilai
policy yang benar-benar dijalankan, termasuk langkah acaknya, sehingga ia memilih jalur yang lebih jauh dari tepi dan
jarang jatuh. Selama eksplorasi masih berjalan, SARSA mendapat ganjaran rata-rata yang lebih baik. Kalau $\epsilon$
diturunkan ke nol, keduanya bertemu di jalur optimal.

Di sini muncul dilema yang tidak dikenal di supervised learning: \istilah{exploration} melawan
\istilah{exploitation}. Kalau agen selalu memilih aksi yang sejauh ini terbaik, ia mungkin tidak pernah
menemukan aksi yang sebenarnya lebih baik. Kalau ia terus mencoba hal baru, ia membuang ganjaran.
Strategi paling sederhana adalah $\epsilon$-greedy: pilih aksi terbaik hampir selalu, tetapi dengan
peluang kecil $\epsilon$ pilih aksi acak. Ini mirip memilih tempat makan. Kalau selalu ke warung
langganan, kita tidak pernah tahu ada warung yang lebih enak di gang sebelah.

\section{Deep Q-network}

Untuk keadaan yang amat banyak, seperti piksel layar permainan, tabel $Q$ diganti jaringan saraf
$Q_\theta(s,a)$. Gabungan ini ternyata tidak stabil, karena data yang dikumpulkan berurutan sangat
berkorelasi dan targetnya terus berubah seiring jaringannya sendiri berubah. DQN \citep{mnih2015}
memperbaikinya dengan dua trik. \istilah{Experience replay} menyimpan pengalaman di sebuah buffer dan
melatih jaringan dengan mini-batch acak dari buffer itu, sehingga korelasi terputus. \istilah{Target
network} adalah salinan jaringan yang diperbarui jarang-jarang dan dipakai untuk menghitung target,
sehingga target tidak bergeser di setiap langkah. Dengan dua trik ini, satu arsitektur yang sama
belajar memainkan puluhan permainan Atari dari piksel. Operator maksimum dalam target Q-learning
cenderung melebih-lebihkan nilai. Alasannya statistik: kalau sepuluh aksi sebenarnya bernilai sama, tetapi perkiraannya
berisik, maksimum dari sepuluh perkiraan berisik hampir selalu lebih besar dari nilai sebenarnya. Double DQN memisahkan
pemilihan aksi dari penilaiannya: jaringan utama memilih aksi terbaik, $a^*=\argmax_{a'}Q_\theta(s',a')$, dan jaringan
target menilai aksi itu, $Q_{\theta^-}(s',a^*)$. Karena noise kedua jaringan tidak sama, bias ke atas itu sebagian besar
hilang.

\section{Policy gradient}

Pendekatan lain melewati fungsi nilai dan langsung mengoptimalkan policy berparameter
$\pi_\theta(a\mid s)$. Teorema policy gradient memberi gradien return harapan $J(\theta)$ dalam bentuk
yang bisa diperkirakan dari sampel,
\begin{equation}
\nabla_\theta J(\theta)=\E_{\pi_\theta}\Big[\sum_t\nabla_\theta\log\pi_\theta(a_t\mid s_t)\,G_t\Big].
\end{equation}
Kuncinya adalah \istilah{log-derivative trick}. Untuk satu lintasan $\tau$ dengan peluang $p_\theta(\tau)$
dan return $G(\tau)$, gradien return harapan adalah
\[
\nabla_\theta\int p_\theta(\tau)G(\tau)\dd\tau=\int p_\theta(\tau)\,\nabla_\theta\log p_\theta(\tau)\,G(\tau)\dd\tau,
\]
karena $\nabla p=p\,\nabla\log p$. Bentuk kedua adalah sebuah ekspektasi, sehingga bisa diperkirakan dengan
merata-ratakan atas lintasan yang dijalankan. Peluang lintasan adalah hasil kali peluang aksi dari
policy dan peluang transisi dari lingkungan. Setelah logaritma diambil, hasil kali itu menjadi jumlah,
dan suku-suku transisi lingkungan tidak bergantung pada $\theta$, sehingga hilang saat diturunkan.
Yang tersisa hanya $\sum_t\nabla_\theta\log\pi_\theta(a_t\mid s_t)$. Inilah alasan policy gradient bisa
dipakai tanpa model lingkungan sama sekali. Mengganti $G(\tau)$ dengan return sejak langkah $t$ saja
diperbolehkan, karena aksi di langkah $t$ tidak bisa memengaruhi ganjaran yang sudah diterima sebelumnya.

Algoritme REINFORCE \citep{williams1992} memakai rumus ini langsung: naikkan peluang aksi yang
diikuti return tinggi, turunkan peluang aksi yang diikuti return rendah. Estimasi ini tak bias,
tetapi variansnya besar. Mengurangkan sebuah \istilah{baseline} yang tidak bergantung pada aksi, biasanya
perkiraan $V(s_t)$, tidak mengubah ekspektasinya tetapi mengurangi variansnya. Selisih
$A_t=G_t-V(s_t)$ disebut \istilah{advantage}: seberapa jauh aksi ini lebih baik daripada rata-rata.
Metode \istilah{actor-critic} melatih dua jaringan sekaligus, aktor yang memilih aksi dan kritikus yang
menaksir nilai \citep{mnih2016}, dan \istilah{generalized advantage estimation} \citep{schulman2016}
mencampur perkiraan advantage dari berbagai panjang langkah untuk mengatur keseimbangan bias dan
varians.

Policy gradient punya kelemahan praktis: satu langkah yang terlalu besar bisa merusak policy, dan
data baru yang dikumpulkan dengan policy rusak itu ikut buruk, sehingga pelatihan bisa runtuh.
TRPO \citep{schulman2015} membatasi setiap update agar policy baru tidak terlalu jauh dari yang lama
dalam arti divergensi KL. PPO \citep{schulman2017} mencapai efek serupa dengan cara yang jauh lebih
sederhana, yaitu memotong rasio peluang $\rho_t=\pi_\theta(a_t\mid s_t)/\pi_{\theta_{\mathrm{lama}}}(a_t\mid s_t)$:
\begin{equation}
L^{\mathrm{CLIP}}(\theta)=\E_t\Big[\min\big(\rho_tA_t,\ \mathrm{clip}(\rho_t,1-\epsilon,1+\epsilon)A_t\big)\Big].
\end{equation}
Kalau sebuah aksi bagus ($A_t>0$), peluangnya boleh naik, tetapi tidak lebih dari faktor $1+\epsilon$
dalam satu update. Kalau aksi buruk, peluangnya boleh turun, tetapi tidak lebih dari faktor
$1-\epsilon$. Contoh dengan angka: dengan $\epsilon=0{,}2$ dan advantage $A_t=2$, kalau policy baru menaikkan peluang aksi itu menjadi
$1{,}5$ kali, suku pertama memberi $3$, tetapi suku yang dipotong memberi $1{,}2\cdot2=2{,}4$, dan minimum keduanya, $2{,}4$, yang
dipakai. Gradien terhadap $\rho_t$ di daerah itu nol, sehingga tidak ada dorongan untuk menaikkan peluangnya lebih jauh dalam
update yang sama. PPO menjadi algoritme bawaan di banyak aplikasi karena kuat dan mudah disetel. Untuk
aksi kontinu, soft actor-critic \citep{haarnoja2018} menambahkan entropi policy ke dalam tujuan,
sehingga agen dihargai karena tetap acak selama hal itu tidak terlalu merugikan.

\section{Meniru dan memberi kredit}

\subsection{Imitation learning}

Kadang lebih mudah menunjukkan cara melakukan sesuatu daripada menulis fungsi ganjarannya.
\istilah{Behavioral cloning} memperlakukan demonstrasi ahli sebagai data supervised: input keadaan,
target aksi. Masalahnya adalah \istilah{covariate shift}. Begitu agen sedikit menyimpang dari jalur
ahli, ia sampai di keadaan yang tidak pernah dilihat ahli, membuat kesalahan lebih besar, dan
menyimpang semakin jauh. DAgger \citep{ross2011} mengatasinya dengan menjalankan policy agen, meminta
ahli memberi label pada keadaan yang dikunjungi agen, lalu melatih ulang. Pendekatan lain
mempelajari fungsi ganjaran dari demonstrasi (\istilah{inverse reinforcement learning}), atau
mempertandingkan agen dengan diskriminator yang membedakan perilakunya dari perilaku ahli, seperti
pada GAIL \citep{ho2016}.

\subsection{Credit assignment}

Masalah tersulit RL adalah ganjaran yang tertunda. Dalam sebuah pertandingan sepak bola, gol terjadi
di menit ke-80, tetapi umpan yang menentukan mungkin terjadi lima menit sebelumnya, dan kesalahan
lawan yang membuka ruang mungkin terjadi lebih awal lagi. Metode TD menyebarkan kredit mundur satu
langkah demi satu langkah, sehingga untuk penundaan panjang, belajarnya lambat sekali. RUDDER
\citep{arjonamedina2019}, yang dikembangkan di Linz, mengambil pendekatan berbeda. Sebuah LSTM dilatih
untuk memprediksi return total dari seluruh urutan keadaan dan aksi. Lalu, dengan menganalisis
seberapa besar setiap langkah mengubah prediksi itu, return dibagi ulang ke langkah-langkah yang
benar-benar berkontribusi. Ganjaran yang tertunda diubah menjadi ganjaran langsung, dan nilai harapan
yang tersisa untuk dipelajari menjadi hampir nol. Kuliah deep RL di Linz membahas topik ini dengan
nama direct credit assignment.

\subsection{Kendali sebagai inferensi}

Kuliah juga membahas cara pandang yang menyatukan RL dengan inferensi probabilistik (\cref{ch:prob}).
Kalau kita memperkenalkan variabel ``optimalitas'' yang peluangnya sebanding dengan eksponensial
ganjaran, mencari policy optimal menjadi mirip menghitung posterior atas lintasan, dengan syarat
bahwa lintasan itu optimal. Hasil dari cara pandang ini adalah RL dengan entropi maksimum, yang
menjelaskan mengapa soft actor-critic menambahkan entropi ke dalam tujuannya.

\section{Model, permainan, dan banyak agen}

Buku \citet{plaat2022}, yang menjadi bacaan kuliah, menyusun deep RL menjadi beberapa keluarga. Setiap keluarga menjawab
satu kelemahan metode model-free yang dibahas di atas.

\subsection{Belajar dengan model}

Metode model-free butuh interaksi yang amat banyak. DQN memainkan setiap permainan Atari puluhan juta langkah. Metode berbasis
model mempelajari dinamika lingkungan, $\hat p(s',r\mid s,a)$, lalu memakai model itu untuk merencanakan atau untuk menghasilkan
pengalaman buatan. Arsitektur Dyna \citep{sutton1990} adalah contoh paling sederhana. Setiap langkah nyata dipakai dua kali:
sekali untuk memperbarui $Q$ seperti biasa, dan sekali untuk memperbarui model. Lalu, sebelum langkah nyata berikutnya, agen
melakukan beberapa pembaruan $Q$ tambahan dengan pengalaman yang dibangkitkan model dari keadaan yang pernah dikunjungi. Kalau
sepuluh langkah khayalan dilakukan untuk setiap langkah nyata, informasi dari satu interaksi menyebar jauh lebih cepat ke
seluruh tabel nilai.

Bahayanya adalah galat model. Policy yang dioptimalkan di dalam model yang keliru akan mengeksploitasi kekeliruan itu, persis
seperti optimasi molekul yang mengeksploitasi celah model QSAR (\cref{ch:hayati}). Galat juga menumpuk sepanjang rencana yang
panjang, masalah rollout yang sama dengan \cref{ch:operator}. Karena itu metode modern biasanya merencanakan hanya beberapa
langkah ke depan, memakai ensemble model untuk menaksir ketidakpastiannya, atau seperti MuZero, mempelajari model yang hanya
meramalkan besaran yang dibutuhkan untuk merencanakan, yaitu nilai, policy, dan ganjaran, bukan seluruh pengamatan.

\subsection{Banyak agen}

Ketika beberapa agen belajar bersamaan, lingkungan setiap agen berubah karena agen lain juga berubah, dan jaminan konvergensi
dari MDP tunggal hilang. Teori permainan memberi bahasa untuk situasi ini. Sebuah \istilah{Nash equilibrium} adalah kombinasi
strategi di mana tidak ada pemain yang bisa untung dengan mengubah strateginya sendiri sementara yang lain tetap.

Batu-gunting-kertas memperlihatkan mengapa strategi acak bisa optimal. Misalkan lawan memainkan batu, gunting, dan kertas
dengan peluang $p_b$, $p_g$, dan $p_k$, dengan menang bernilai $+1$, kalah $-1$, dan seri $0$. Nilai harapan memainkan batu
adalah $p_g-p_k$, gunting $p_k-p_b$, dan kertas $p_b-p_g$. Kalau salah satu nilai ini positif, kita akan selalu memainkan aksi
itu, dan lawan bisa mengeksploitasinya. Satu-satunya keadaan di mana tidak ada aksi yang lebih unggul adalah ketika ketiganya
nol, yaitu $p_b=p_g=p_k=\tfrac13$. Ekuilibriumnya adalah bermain acak sempurna. Dalam permainan seperti ini, self-play yang
naif bisa berputar-putar tanpa henti: agen belajar mengalahkan versi lamanya, lalu versi berikutnya belajar mengalahkan versi
itu, dan seterusnya. Karena itu sistem seperti AlphaZero \citep{silver2018} dan agen permainan strategi modern melatih agen
melawan kumpulan versi-versi lamanya, bukan hanya versi terakhir.

\subsection{Hierarki, meta-learning, dan eksplorasi}

Tugas yang panjang, seperti memasak sebuah resep, lebih mudah dipelajari kalau dipecah menjadi subtugas: ambil bahan, potong,
panaskan. Kerangka \istilah{options} \citep{sutton1999} memperlakukan subtugas seperti itu sebagai aksi yang diperluas dalam
waktu, masing-masing dengan policy dan syarat berhentinya sendiri. Policy tingkat tinggi memilih option, policy tingkat rendah
menjalankannya. Credit assignment menjadi lebih mudah, karena keputusan tingkat tinggi lebih sedikit dan berjarak lebih dekat
dengan ganjarannya.

Meta-learning bertanya bagaimana agen bisa cepat beradaptasi ke tugas baru. MAML \citep{finn2017} mencari parameter awal $\theta$
yang, setelah satu atau beberapa langkah gradien pada tugas baru, sudah bekerja baik di tugas itu. Untuk setiap tugas $i$,
parameter yang sudah disesuaikan adalah $\theta_i'=\theta-\alpha\nabla_\theta\Loss_i(\theta)$, dan parameter awal diperbarui untuk
meminimalkan $\sum_i\Loss_i(\theta_i')$, yaitu loss \emph{setelah} adaptasi. Yang dicari bukan parameter yang baik untuk semua
tugas sekaligus, melainkan titik awal yang dekat ke solusi setiap tugas. Gagasan ini berlaku di luar RL, misalnya untuk
klasifikasi dengan sedikit contoh.

Eksplorasi yang cerdas menjadi penting ketika ganjaran jarang sekali. Bonus eksplorasi berbasis hitungan menambahkan ganjaran
buatan $\beta/\sqrt{N(s)}$ untuk keadaan yang baru dikunjungi $N(s)$ kali, gagasan yang sama dengan bonus keingintahuan pada UCT
di \cref{ch:klasik}. Untuk keadaan berupa piksel, setiap layar praktis unik, sehingga hitungan langsung tidak berguna.
\citet{bellemare2016} menurunkan ``hitungan semu'' dari sebuah model kerapatan atas keadaan, dan memperlihatkan bahwa bonus
seperti ini membuat agen menjelajahi permainan Atari yang ganjarannya jarang jauh lebih baik.

\section{Reinforcement learning di dunia nyata}

Di luar permainan, RL dipakai untuk masalah kendali yang sulit dimodelkan secara analitis.
\citet{rabault2019} melatih agen untuk mengendalikan jet kecil di permukaan silinder dalam simulasi
CFD, dan agen itu menemukan strategi yang mengurangi gaya hambat. Ulasan \citet{garnier2021}
merangkum puluhan studi serupa dalam mekanika fluida, dari kendali aliran sampai optimasi bentuk.
Dalam robotika, agen dilatih di simulator lalu dipindahkan ke robot nyata. Dalam sistem rekomendasi,
memilih item yang ditampilkan kepada pengguna bisa dirumuskan sebagai masalah keputusan berurutan.

Penerapan yang paling berpengaruh beberapa tahun terakhir justru datang dari bahasa. Model bahasa
yang sudah dilatih kemudian disetel dengan \istilah{reinforcement learning from human feedback}.
Manusia membandingkan pasangan jawaban model, sebuah model ganjaran dilatih dari preferensi itu
\citep{christiano2017}, lalu model bahasa dioptimasi dengan PPO terhadap model ganjaran tersebut
\citep{ouyang2022}. Perkembangan ini kita bahas lagi di \cref{ch:terbaru}.

\section{Perkembangan riset}

Dua arah riset yang menonjol beberapa tahun terakhir menyentuh dua pertanyaan besar bab ini: dari mana data datang, dan
bagaimana memakai model. DreamerV3 \citep{hafner2023} mempelajari model dunia dari pengalaman, lalu memperbaiki
perilakunya dengan membayangkan skenario masa depan di dalam model itu. Dengan teknik normalisasi dan transformasi yang
membuat skala ganjaran dan nilai tidak lagi perlu disetel per tugas, satu konfigurasi yang sama bekerja di lebih dari 150
tugas, dan untuk pertama kalinya sebuah algoritme mengumpulkan berlian di Minecraft dari nol tanpa data manusia atau
kurikulum.

Decision Transformer \citep{chen2021dt} membuang fungsi nilai dan policy gradient sama sekali. Untuk data yang sudah
terkumpul (offline RL), setiap lintasan ditulis sebagai deret token berisi return yang tersisa, keadaan, dan aksi.
Transformer dengan causal mask dilatih menebak aksi berikutnya. Saat dipakai, kita cukup memasukkan return yang
diinginkan sebagai token pertama, dan model membangkitkan aksi yang konsisten dengan target itu. Reinforcement learning
berubah menjadi pemodelan deret bersyarat, dan seluruh perkakas model bahasa dari \cref{ch:lstm} bisa dipakai. Pendekatan
ini juga menjadi jembatan ke cara model bahasa besar dilatih dengan RL di \cref{ch:terbaru}.

\sumberbab{Bab ini menggabungkan kuliah Reinforcement Learning (MDP, dynamic programming, Monte Carlo,
temporal difference, policy gradient, PPO) dan kuliah serta latihan Deep Reinforcement Learning
(imitation learning, policy gradient, direct credit assignment, kendali dan inferensi, proyek DQN dan
PPO). Bacaan utama: \citet{sutton2018} dan \citet{plaat2022}.}
